\subsection{Momentum and thermal closures}
\label{subsec:wall-functions:gao-li:governing-relations}

The thermal closure is a quadratic relation between static enthalpy and velocity~\cite[Eqs.~(12)--(15)]{li2025WF}. It extends the generalized temperature--velocity relation of Zhang et al. by using enthalpy to accommodate reacting-gas thermodynamics. Its generalized recovery enthalpy determines the wall slope and thereby relates heat-flux to wall shear, extending the fixed-recovery-factor prescription of the earlier Walz-based construction.
\par\medskip

The momentum closure combines a viscosity-corrected viscous-sublayer relation with a compressible logarithmic law in a unified expression~\cite[Eqs.~(27), (30)--(32)]{li2025WF}. It assumes approximately constant total shear stress in the inner layer and blends through the buffer layer. Sampling below the logarithmic region is therefore permitted by construction, although accuracy can still vary with sampling height.
\par\medskip

\label{subsec:wall-functions:gao-li:empirical-constants}
The formulation retains \(\kappa=0.41\) and \(B=6.0\) and approximates the viscosity variation by an enthalpy power law. Li et al. fit its exponent, denoted \(\omega_H\) here, to approximately \(0.72\) in their reacting viscous-sublayer profiles~\cite[Eq.~(25) and Figure~12]{li2025WF}. This transport approximation is calibrated for those conditions; its transfer to other gas models and wall temperatures is not automatically given and open to scrutiny.
